
%%%===========================================
\section{Organization of the Thesis}
% -- Sugestão: Esta subseção descreve a estrutura do trabalho, ajudando o leitor a entender a progressão das seções. Ajuste conforme o desenvolvimento da tese, mas mantenha a ideia de profundidade crescente: Introduction apresenta o problema e objetivos, Literature Review analisa a literatura, Theoretical Framework propõe o economic commensalism, Mathematical Modeling formaliza o modelo, Empirical Validation testa empiricamente, e Conclusion sumariza os achados e implicações.

Following this general hypothesis, the study is organized into three complementary chapters, each addressing a specific dimension of the problem — theoretical, formal, and empirical — through the articulation of specific hypotheses:

\begin{itemize}
    \item \textbf{H1 (Theoretical - Chapter 1):} Incorporating the ecological concept of commensalism into economic theory enables a more coherent and incentive-compatible framework for aligning productive rationality with environmental preservation.
    
    \item \textbf{H2 (Mathematical - Chapter 2):} In a dynamic optimization model incorporating Bellman recursion, Kalman filtering, Bayesian learning, and agent interactions (Spin), it is possible to derive a sustainable Nash equilibrium based on incentive-compatible environmental policies.
    
    \item \textbf{H3 (Empirical - Chapter 3):} In the context of rural producers in Mato Grosso, incentive-based environmental policies inspired by commensal practices are more effective at promoting conservation behavior than punitive or diffuse compensation mechanisms.
\end{itemize}

Each of these specific hypotheses supports and refines the general hypothesis, forming a coherent framework in which theoretical insights, formal logic, and empirical evidence interact to validate (or falsify) the proposed paradigm.


This thesis is organized as follows. Section 2 reviews the literature on environmental economics, focusing on the tragedy of the commons, market-based solutions, game theory, and behavioral economics, identifying key gaps that the proposed framework addresses. Section 3, the Theoretical Framework, builds on this review to formally define \textit{economic commensalism} and propose hypotheses for reconciling economic incentives with environmental preservation. Section 4, Mathematical Modeling, formalizes the framework using dynamic optimization, uncertainty modeling, and game-theoretic simulations. Section 5, Empirical Validation and Discussion, tests the model through econometric analysis and case studies in Brazil's Mato Grosso region, discussing the results in relation to the literature and objectives. Section 6 concludes with theoretical and practical contributions, policy implications, limitations, and directions for future research. Appendices include a timeline and a note on the use of AI.

%%%===========================================
\subsection{Theoretical Framework}
% -- Sugestão: Este capítulo deve ser o primeiro dos três principais, como você planejou. Use a revisão de literatura para sintetizar as lacunas identificadas e propor o economic commensalism como solução. Defina o conceito com rigor (e.g., um mecanismo de incentivos que alinha interesses privados com sustentabilidade, respeitando limites biofísicos), formalize-o teoricamente (e.g., com base em teoria dos jogos ou economia comportamental), e derive hipóteses testáveis que serão formalizadas no Mathematical Modeling e testadas na Empirical Validation.
% -- Vai responder os gaps da literature review
% -- Vai definir os parâmetros teóricos a serem testados de forma dedutiva (cap. 2) e indutiva (cap. 3).

\subsubsection{Conceptual Development}
\begin{itemize}
    \item \textbf{a:} aaa
    \item \textbf{b:} bbb
    \item \textbf{c:} ccc
\end{itemize}

\subsubsection{Proposed Framework: Economic Commensalism}
% Sugestão: Introduza o economic commensalism com uma definição clara, mostrando como ele resolve as lacunas identificadas na Literature Review. Por exemplo, "Economic commensalism is a market-based mechanism that uses rational incentives to align private interests with collective environmental goals, addressing the scalability issues of Ostrom’s self-governance and the uncertainty limitations of Pigou’s taxes." Derive hipóteses específicas que serão testadas nos capítulos seguintes.

[Conteúdo a ser preenchido pelo autor.]

\subsubsection{Hypotheses Formulation}
% -- Sugestão: Refine as hipóteses da Introduction para serem mais específicas ao economic commensalism, garantindo que sejam testáveis no Mathematical Modeling e Empirical Validation. Por exemplo, "Under the economic commensalism framework, well-structured incentives will shift producers from an exploitation Nash equilibrium to a preservation equilibrium."

[Conteúdo a ser preenchido pelo autor, com base nas hipóteses da Introduction.]

%%%===========================================
\subsection{Mathematical Modeling}
% ver  (C204052025a)
% -- Sugestão: Este capítulo deve formalizar o economic commensalism em um modelo matemático, como você planejou. Use as hipóteses do Theoretical Framework para guiar a modelagem. Por exemplo, se uma hipótese é que incentivos podem mudar o equilíbrio de Nash, desenvolva um modelo de jogo que incorpore incentivos e externalidades. Certifique-se de que o modelo gere previsões testáveis (e.g., impacto de incentivos na redução do desmatamento), que serão avaliadas na Empirical Validation.

\subsubsection{Dynamic Optimization}
Bellman’s equation for intertemporal decision-making.

\subsubsection{Uncertainty Modeling}
Kalman filters for adjusting predictions based on real-time information.

\subsubsection{Agent-Based Learning}
Bayesian-Markovian models to capture adaptive expectations.

\subsubsection{Game-Theoretic Simulations}
Pareto and Nash equilibrium to assess policy effectiveness.

%%%===========================================
\subsection{Empirical Validation and Discussion}
% -- Sugestão: Este capítulo combina a validação empírica com a discussão dos resultados, como sugerido. Aplique o modelo matemático ao Mato Grosso, usando dados específicos (e.g., taxas de desmatamento, custos de produção) e métodos claros (e.g., regressão, análise de impacto). Teste as hipóteses do Theoretical Framework e discuta os resultados em relação à literatura (e.g., Pigou, Ostrom) e aos objetivos da tese. Mostre como os achados validam (ou não) o economic commensalism.

Cap. 1: lacunas teóricas e proposições para solucioná-las
Cap. 2: formalização dedutiva e coerência interna
Cap. 3: validação empírica e coerência externa

\subsubsection{Empirical Validation}
\begin{itemize}
    \item Analysis of historical environmental policy outcomes (California Cap-and-Trade, REDD+ initiatives).
    \item Field surveys with agricultural producers in Mato Grosso to assess incentive responsiveness.
    \item Econometric testing of incentive efficiency using panel data regression and machine learning techniques.
\end{itemize}

\subsubsection{Discussion}
% -- Sugestão: Compare os resultados com a literatura (e.g., como o economic commensalism supera as limitações de Pigou ou Ostrom) e avalie a viabilidade prática da proposta. Discuta se os achados apoiam as hipóteses e como eles contribuem para os objetivos da tese. Identifique limitações do estudo (e.g., dados disponíveis, escalabilidade) para preparar a Conclusion.

[Conteúdo a ser preenchido pelo autor.]

\newpage

% fim %